% !TEX root = report.tex
% !TEX program = pdflatex

\documentclass[11pt]{article}

\usepackage[a4paper,margin=1in]{geometry}
\usepackage{amsmath,amssymb}
\usepackage{booktabs}
\usepackage{siunitx}
\usepackage{hyperref}

\title{Guaranteeing Physical Validity by Construction: A Systematic
Factorial Evaluation of Hard vs.\ Soft Thermodynamic Constraint
Enforcement in Physics-Informed Neural Networks for Sorption Studies}
\author{}
\date{}

\begin{document}
\maketitle

\begin{abstract}
    To be written after the main results and discussion have been finalized.
\end{abstract}

\section{Introduction}

To be written.

\section{Materials and Methods}

\subsection{Dataset and prediction task}

The study considers methane--coal sorption data under different pressure,
temperature, and coal-composition conditions. The raw dataset contains
993 adsorption measurements. Because the CSV did not include experiment
identifiers, experiment groups were constructed from rows with the same coal
composition and temperature, resulting in 114 experiment groups. Rows in the
same experiment group are treated as one adsorption curve, meaning a set of
pressure--adsorption measurements collected at fixed coal composition and
temperature. Splitting was performed at the experiment level to avoid placing
points from the same curve in both training and test sets.

The raw input columns are moisture content, ash content, volatile matter,
temperature, and pressure. The prediction target is adsorption capacity. The
corresponding variables used by the model are summarized as follows:
\begin{center}
    \begin{tabular}{ll}
        \toprule
        \textbf{Role}         & \textbf{Model variable}       \\
        \midrule
        Temperature input     & \texttt{temperature\_K}       \\
        Pressure input        & \texttt{pressure\_MPa}        \\
        Moisture input        & \texttt{moisture\_wt}         \\
        Ash input             & \texttt{ash\_wt}              \\
        Volatile-matter input & \texttt{volatile\_wt}         \\
        Prediction target     & \texttt{CH4\_adsorption\_m3t} \\
        \bottomrule
    \end{tabular}
\end{center}

Since thermodynamic formulas require absolute temperature, the measured values
were converted from degrees Celsius to Kelvin before any thermodynamic feature
or constraint was evaluated. Since the Sips formulation requires raw positive
pressure, all samples were checked to satisfy \(P>0\). Although adsorption
capacity is physically non-negative, the dataset was checked to ensure strictly
positive target values, since the logarithmic target transform requires
\(q>0\).

The raw CSV did not include a measured \texttt{coal\_rank} column. To define
the T2 coal-rank extrapolation tier, coal-rank labels were constructed from the
available coal-composition measurements: moisture, ash, and volatile matter.
Volatile matter and fixed carbon were first expressed on a dry, ash-free basis,
and threshold rules were then used to assign each sample to one coal rank. This
produced four ranks in the analyzed dataset: anthracite, low-volatile
bituminous, medium-volatile bituminous, and subbituminous/lignite. The
constructed coal-rank label is used only for T2 grouping and interpretation
analyses, not as a model input.

\subsection{Preprocessing and feature engineering}

Preprocessing makes the data valid, comparable, and suitable for training.
Engineered features were also included to represent derived thermodynamic and
coal-composition information. The model input vector contains 12 numerical
features:
\[
    \mathbf{x} =
    [
    \mathrm{moisture},
    \mathrm{ash},
    \mathrm{volatile},
    T,
    P,
    T_r,
    P_r,
    \mathrm{FC},
    \mathrm{OM},
    \beta,
    PT,
    \mathrm{MV}
    ].
\]
Here,
\[
    \begin{aligned}
        T_r         & = T/190.6,                                              \\
        P_r         & = P/4.60,                                               \\
        \mathrm{FC} & = 100-\mathrm{moisture}-\mathrm{ash}-\mathrm{volatile}, \\
        \mathrm{OM} & = (\mathrm{volatile}+\mathrm{FC})/100,                  \\
        \beta       & = (RT)^{-1},                                            \\
        PT          & = P T,                                                  \\
        \mathrm{MV} & = \mathrm{moisture}\times\mathrm{volatile}.
    \end{aligned}
\]
The gas constant was set to \(R=8.314~\mathrm{J\,mol^{-1}\,K^{-1}}\).


All input features were scaled using only the training data, so that the model
could train on variables with comparable numerical ranges without using
information from the validation or test sets. The adsorption target was
log-transformed during training, and predictions were converted back to the
original adsorption scale before computing the reported error metrics. Raw
pressure in MPa was also kept separately, because the Sips-based model
components and physical checks require the real pressure values rather than the
scaled pressure feature.

\subsection{Leakage-safe data splitting}

To avoid data leakage, all splits were performed at the experiment-group level
rather than at the individual-row level. Rows belonging to the same adsorption
curve share the same \texttt{experiment\_id}, and were therefore kept together
in the same split. This prevents the model from being trained on part of an
adsorption curve and then tested on another pressure point from the same curve.

The data pipeline defines three evaluation tiers. T1 is the interpolation test
set, containing held-out experiment groups from coal ranks that are also present
in training. T2 is the coal-rank extrapolation set, where one coal-rank group is
withheld from training. T3 is the pressure extrapolation set, defined by
high-pressure samples inside the T1 test set. Cross-validation was also
performed with GroupKFold using \texttt{experiment\_id} as the grouping
variable, so that validation folds remain leakage-safe.

\subsection{Soft-constrained baseline}
The soft-constrained baseline was implemented as the reference model for the
study. It uses the 12 engineered input features to predict methane adsorption
capacity, while also estimating auxiliary physical quantities used in the
physics penalties. Physical consistency is encouraged through additional loss
terms for saturation bounds, pressure monotonicity, Sips consistency, and
temperature dependence, but these constraints are not guaranteed by the model
architecture.

A sweep over the physics-loss weight \(\lambda_p\) was performed
to examine the trade-off between predictive accuracy and physical consistency.
The trained baseline was evaluated on the T1 interpolation, T2 coal-rank
extrapolation, and T3 pressure-extrapolation tiers, and serves as the reference
for comparison with the hard-constrained variants.


\subsection{Hard-constrained model variants}

Hard-constrained model variants were implemented to compare against the
soft-constrained baseline and to test whether direct architectural enforcement
improves physical validity and extrapolation performance. In this project, the
hard variants are:
\[
    \begin{array}{ll}
        \mathrm{V3}    & \text{hard bounds},                             \\
        \mathrm{V4}    & \text{hard pressure monotonicity},              \\
        \mathrm{V5}    & \text{hard Sips gauge},                         \\
        \mathrm{V6}    & \text{hard Van't Hoff/Arrhenius structure},     \\
        \mathrm{V7}    & \text{joint hard bounds and hard monotonicity}, \\
        \mathrm{V2/V8} & \text{all-hard combined variants}.
    \end{array}
\]

The hard-bounds variant enforces the adsorption range directly,
\[
    0 < \hat{q} < \hat{q}_{\max}.
\]
Instead of penalizing saturation violations after they occur, the output head
is designed so that the adsorption prediction remains between zero and the
learned maximum capacity.

The hard-monotonicity variant enforces non-decreasing adsorption with respect to pressure,
\[
    \frac{\partial \hat{q}}{\partial P} \geq 0.
\]
Pressure enters through a monotone pathway with non-negative weights and
monotone activations. This makes the pressure contribution unable to decrease
the predicted adsorption. In the implementation, raw pressure and the
pressure-derived engineered features \(P_r\) and \(PT\) were excluded from the
unconstrained composition branch. This is necessary because an unconstrained
branch could otherwise learn a decreasing pressure effect and break the
monotonicity guarantee.

The hard Sips-gauge variant does not predict adsorption as a completely free
output. Instead, it predicts the Sips parameters
\[
    \hat{q}_{\max}, \quad \hat{K}, \quad \hat{n},
\]
from coal-composition features, and then computes adsorption using raw pressure:
\[
    \hat{q}(P)
    =
    \hat{q}_{\max}
    \frac{(\hat{K}P)^{\hat{n}}}
    {1+(\hat{K}P)^{\hat{n}}}.
\]
This gives the model a structured pressure--adsorption relation.

The hard Van't Hoff/Arrhenius variant constrains the temperature dependence of
the adsorption constant. The model predicts an Arrhenius-type structure so that
the relationship between \(K\), temperature, and adsorption enthalpy \(\Delta H\) is built
into the model. Finally, the joint hard variants combine multiple hard constraints to test
whether the physical rules work well together or whether they reduce accuracy
or extrapolation performance.

\subsection{Training protocol}

All model variants were trained using the same general training protocol. The
training data were provided in mini-batches with a default batch size of 64.
For each batch, the model predicted methane adsorption capacity and the variant-specific
auxiliary physical quantities. The total loss was then computed from the
data-fitting term, together with the physics-consistency penalties for
model variants that used soft penalty terms.

Model parameters were optimized using the AdamW optimizer with learning rate
\(10^{-3}\) and weight decay \(10^{-5}\). The full V1 physics-weight sweep used
the blueprint-specified network architecture with hidden dimensions
\((512,1024,512,256,128)\), a maximum of 1000 epochs, early stopping patience of
100 epochs, three random seeds, and five grouped cross-validation folds. The
physics-loss weights were
\(\lambda_p\in\{0,0.01,0.05,0.1,0.2,0.5,1.0,2.0\}\), giving 15 validation runs
per value of \(\lambda_p\).

Gradient clipping with maximum norm 5 was used to improve numerical stability.

During training, the validation set was used only for model selection and early
stopping. The model state with the lowest validation loss was restored after
training and saved as a checkpoint. Final evaluation was then performed on the
held-out T1, T2, and T3 evaluation tiers using the unified evaluation harness.



\subsection{Evaluation metrics and statistical tests}

All trained models were evaluated using the same unified evaluation harness and
the same T1, T2, and T3 evaluation tiers. Predictive accuracy was measured on
the original adsorption scale after reversing the logarithmic target transform.
The main accuracy metrics were the coefficient of determination \(R^2\), root
mean squared error (RMSE), mean absolute error (MAE), maximum absolute error
(MaxAE), and relative error.

Physical consistency was evaluated using violation rates. These included the
fraction of negative adsorption predictions, the fraction of predictions above
the predicted saturation capacity \(q_{\max}\), and the fraction of
pressure--adsorption monotonicity violations. These metrics were used to assess
whether a model only fit the data numerically or also respected the intended
physical constraints.

Uncertainty quality was evaluated using Monte Carlo dropout predictions and
calibration metrics, including empirical interval coverage and expected
calibration error (ECE). Confidence intervals for accuracy and physical
violation metrics were estimated by bootstrap resampling at the experiment
level, preserving the grouped structure of the data.

For model comparisons, paired statistical tests were performed at the
experiment level. In particular, pairwise RMSE differences between model
variants were computed on the same evaluation samples, and one-sample tests
were applied to the experiment-wise differences. When multiple pairwise
comparisons were made, Bonferroni correction was applied to control for
multiple testing.



\subsection{Uncertainty calibration}

Uncertainty calibration was used to assess whether the predictive uncertainty
estimates were reliable. For each trained model, Monte Carlo dropout was used to
generate repeated stochastic predictions for the same input samples. The spread
of these predictions was used as an estimate of epistemic uncertainty, while
the model-predicted variance represented aleatoric uncertainty when available.

Calibration was evaluated by comparing nominal prediction intervals with their
empirical coverage. For example, a well-calibrated 90\% prediction interval
should contain the observed adsorption value for approximately 90\% of the
samples. Expected calibration error (ECE) was computed as the average difference
between nominal and empirical coverage across the evaluated prediction-interval
levels.

A scalar temperature-scaling factor was selected on the validation set to adjust
the predicted uncertainty widths. The calibrated uncertainty estimates were then
evaluated on the T1, T2, and T3 tiers. This allowed us to test whether the
models became more uncertain under extrapolation, especially for the
coal-rank-extrapolation and high-pressure-extrapolation settings.



\subsection{Interpretability and physical validation}

Interpretability and physical validation were used to examine whether the
trained models learned physically interpretable relationships between pressure,
coal composition, and adsorption behavior. This included feature importance
patterns, Sips parameter trends, and adsorption enthalpy
behavior across coal ranks. Feature importance was evaluated using SHAP-style
analyses for the baseline and Sips-gauge models. The learned Sips parameters
\(\hat{q}_{\max}\), \(\hat{K}\), and \(\hat{n}\) were summarized by coal rank
and compared with classical Langmuir fits. The learned adsorption enthalpy
\(\Delta H\) from the Arrhenius variant was also analyzed by coal rank to check
whether the inferred temperature dependence was physically plausible.



\section{Results and Discussion}

All numerical results reported in this section were obtained from the full
training and evaluation protocol described above.

\subsection{Evaluation setup and data splits}

The processed dataset contained 993 adsorption measurements grouped into 114
constructed experiment groups. The leakage-safe splitting procedure assigned
complete experiment groups to each split, so no constructed experiment
identifier was shared between training, validation, and held-out evaluation data. After withholding
the T2 coal-rank extrapolation set, the remaining data were split into
training, validation, and T1 interpolation test sets. The final split
sizes are shown in Table~\ref{tab:data-splits}.

\begin{table}[htbp]
    \centering
    \begin{tabular}{lrr}
        \toprule
        \textbf{Split}             & \textbf{Samples} & \textbf{Experiments} \\
        \midrule
        Full dataset               & 993              & 114                  \\
        Training pool excluding T2 & 752              & 86                   \\
        Training                   & 472              & 54                   \\
        Validation                 & 132              & 14                   \\
        T1 interpolation           & 148              & 18                   \\
        T2 coal-rank extrapolation & 241              & 28                   \\
        T3 pressure extrapolation  & 3                & 2                    \\
        \bottomrule
    \end{tabular}
    \caption{Leakage-safe data splits used for training, validation, and
        held-out evaluation.}
    \label{tab:data-splits}
\end{table}

The T2 tier consisted of all samples inferred as subbituminous/lignite, giving
a held-out coal-rank extrapolation test with 241 samples from 28 experiments.
The T3 pressure-extrapolation tier contained samples from the T1 set with
pressure above \(8~\mathrm{MPa}\). Because only three samples satisfied this
condition, T3 should be interpreted with caution. Across the full dataset, the
pressure range was
\(0.0019342\)--\(9.2587632~\mathrm{MPa}\), and all pressure values were
strictly positive.

\subsection{Soft baseline and physics-loss weighting}

The soft-constrained baseline was designed to predict methane adsorption while
encouraging physical consistency through additional penalty terms. The total
training objective combines a data-fitting loss with a physics loss,
\[
    \mathcal{L}_{\mathrm{total}}
    =
    \mathcal{L}_{\mathrm{data}}
    +
    \lambda_p \mathcal{L}_{\mathrm{physics}},
\]
where \(\lambda_p\) controls the strength of the physics regularization. A
larger value of \(\lambda_p\) gives more importance to the physical constraints,
whereas \(\lambda_p=0\) corresponds to a purely data-driven baseline.

The implemented physics loss includes terms related to adsorption bounds,
pressure monotonicity, Sips consistency, and temperature dependence. The purpose
of the \(\lambda_p\) sweep is therefore to identify a balance between numerical
accuracy and physical consistency. In practical terms, this analysis asks
whether the model can fit the measured adsorption values while also reducing
violations such as negative adsorption, saturation-bound violations, or
non-monotone pressure response.

The full V1 sweep is shown in Table~\ref{tab:v1-lambda-sweep}. Model selection
used a validation score that combines accuracy and physical validity,
\[
    s_{\mathrm{val}}
    =
    \mathrm{RMSE}_{\mathrm{val}}
    \left(1+v_{\mathrm{val}}\right),
\]
where \(v_{\mathrm{val}}\) is the largest validation violation rate among
negative adsorption, saturation-bound, and non-monotonicity checks. The
unregularized model \((\lambda_p=0)\) had a competitive RMSE but a violation
rate of 1.0. The selected value was \(\lambda_p=0.05\), which gave the lowest
combined validation score while keeping the violation rate close to zero.

\begin{table}[htbp]
    \centering
    \small
    \begin{tabular}{rrrrrr}
        \toprule
        \(\boldsymbol{\lambda_p}\) & \textbf{Val. score}
                                   & \textbf{Val. RMSE}   & \textbf{Val. violation} & \textbf{Val. ECE}
                                   & \textbf{Mean epochs}                                                               \\
        \midrule
        0.00                       & 10.349               & 5.174                   & 1.000             & 0.034 & 116.1 \\
        0.01                       & 5.041                & 5.018                   & 0.004             & 0.049 & 154.5 \\
        0.05                       & \textbf{4.916}       & \textbf{4.897}          & 0.004             & 0.050 & 181.6 \\
        0.10                       & 5.050                & 5.002                   & 0.010             & 0.052 & 213.4 \\
        0.20                       & 5.036                & 5.004                   & 0.006             & 0.055 & 243.9 \\
        0.50                       & 4.969                & 4.966                   & \textbf{0.001}    & 0.056 & 315.3 \\
        1.00                       & 4.941                & 4.922                   & 0.004             & 0.049 & 428.6 \\
        2.00                       & 5.843                & 5.840                   & 0.001             & 0.064 & 445.5 \\
        \bottomrule
    \end{tabular}
    \caption{Full V1 physics-weight sweep using three seeds and five grouped
        cross-validation folds. The selected value is \(\lambda_p=0.05\).}
    \label{tab:v1-lambda-sweep}
\end{table}

\subsection{Comparison of soft and hard variants}

The soft baseline and hard-constrained variants were compared across the three
evaluation tiers using the same evaluation harness.
Table~\ref{tab:rmse-comparison} summarizes the full-protocol RMSE values.
Lower RMSE indicates smaller prediction error on the original methane
adsorption scale.

\begin{table}[htbp]
    \centering
    \begin{tabular}{lrrr}
        \toprule
        \textbf{Variant}               & \textbf{T1 RMSE} & \textbf{T2 RMSE} & \textbf{T3 RMSE} \\
        \midrule
        V1 soft baseline               & 5.300            & 5.746            & 4.943            \\
        V3 hard bounds                 & 5.624            & 5.903            & 2.217            \\
        V4 hard monotonicity           & \textbf{5.213}   & 5.709            & 4.643            \\
        V5 hard Sips gauge             & 5.410            & 5.476            & 1.372            \\
        V6 hard Van't Hoff / Arrhenius & 6.082            & 5.543            & 5.478            \\
        V7 hard bounds + monotonicity  & 7.116            & 6.145            & \textbf{1.192}   \\
        V2/V8 all-hard combined        & 6.209            & \textbf{5.056}   & 1.435            \\
        \bottomrule
    \end{tabular}
    \caption{Full-protocol RMSE values for the soft baseline and
        hard-constrained variants on the T1, T2, and T3 evaluation tiers.}
    \label{tab:rmse-comparison}
\end{table}

For the full-protocol results, the hard monotonicity variant gives the lowest T1
RMSE, although its advantage over the soft baseline is small. On the T2
coal-rank extrapolation tier, the V2/V8 all-hard combined variants give the
lowest RMSE. Thus, in the full-protocol results, combining the Sips-gauge and
Van't Hoff / Arrhenius structures gives the smallest error when the model is
evaluated on a coal rank withheld from training. On the T3
pressure-extrapolation tier, the joint hard-bounds and monotonicity variant gives the
lowest RMSE, but this T3 result is not treated as strong statistical evidence
because T3 contains only three samples.

These results illustrate the main comparison that the hard-variant analysis is
designed to answer: different physical constraints have different effects across
evaluation regimes. A constraint that improves interpolation does not
necessarily give the best extrapolation result, and combining constraints can
improve one tier while reducing performance on another. The reported RMSE
minima differ by evaluation tier: hard constraints help most clearly for
physical validity and for selected extrapolation tiers, but they do not dominate
every accuracy metric.

\subsection{Physical validity of constrained architectures}

Physical validity was evaluated separately from predictive accuracy. This
separation is necessary because a model can have a low numerical error while
still producing negative adsorption, saturation violations, or decreasing
adsorption with increasing pressure. The main physical checks were non-negative
adsorption, adsorption below the predicted saturation capacity, non-decreasing
adsorption with pressure, preservation of the Sips structure, and consistency
with the Arrhenius/Van't Hoff temperature relation.

\begin{table}[htbp]
    \centering
    \begin{tabular}{@{}>{\raggedright\arraybackslash}p{0.23\linewidth}
        >{\raggedright\arraybackslash}p{0.28\linewidth}
        >{\raggedright\arraybackslash}p{0.39\linewidth}@{}}
        \toprule
        \textbf{Variant} & \textbf{Physical check}                         & \textbf{Verification result} \\
        \midrule
        V4 hard monotonicity
                         & \(\partial \hat{q}/\partial P \geq 0\)
                         & minimum slope \(=0.111\), violation rate \(=0\)                                \\
        V5 hard Sips gauge
                         & Sips bounds and monotonicity
                         & lower, upper, and slope violation rates \(=0\)                                 \\
        V6 hard Van't Hoff / Arrhenius
                         & Arrhenius identity
                         & maximum identity error \(=1.82\times10^{-12}\)                                 \\
        V7 hard bounds + monotonicity
                         & \(\partial \hat{q}/\partial P \geq 0\)
                         & minimum slope \(=0.146\), violation rate \(=0\)                                \\
        V2/V8 all-hard combined
                         & Sips bounds and monotonicity
                         & lower, upper, and slope violation rates \(=0\)                                 \\
        \bottomrule
    \end{tabular}
    \caption{Architecture-level physical verification checks for the
        full-protocol hard-constrained variants.}
    \label{tab:physical-verification}
\end{table}

The hard monotonicity variants passed the pressure-monotonicity verification:
the computed pressure derivatives were non-negative on the tested pressure
grid. The hard Sips-gauge variants also satisfied the Sips-bounded,
pressure-monotone structure.
For V5, the minimum predicted adsorption was positive
\((\min\hat{q}=0.0035)\), the minimum margin
\(\hat{q}_{\max}-\hat{q}\) was positive \((0.029)\), and no lower-bound,
upper-bound, or negative-slope violations were detected. This confirms that the
Sips-gauge architecture enforces the desired pressure--adsorption shape by
construction, rather than relying only on a penalty term during training.

The Arrhenius variants satisfied the algebraic Arrhenius identity to numerical
precision. However, the learned adsorption enthalpy values require more careful
interpretation. For V6, the mean learned adsorption enthalpy was approximately
\(-44.2~\mathrm{kJ\,mol^{-1}}\), and about 76\% of sampled values lay in the
expected range used by the verification script. For the V2/V8 all-hard combined
variant, the
Arrhenius identity was also satisfied, but the sampled enthalpy values did not
fall in that expected range. Thus, satisfying the mathematical form of a
constraint does not automatically guarantee that all learned physical
parameters are scientifically plausible.

Overall, the physical-validity checks show the difference between soft and hard
constraints. Soft constraints discourage violations through the loss function,
whereas hard constraints remove certain violations from the model class itself.
The full-protocol verification confirms that the implemented hard architectures
enforce the targeted bounds, monotonicity, Sips-form, and Arrhenius-form
properties.

\subsection{Ablation and interaction effects}

The ablation analysis was used to separate the effect of individual hard
constraints from the effect of combining them. Two factorial comparisons were
implemented. The first comparison studies hard bounds and hard monotonicity
using V1, V3, V4, and V7. The second comparison studies the Sips-gauge and
Van't Hoff / Arrhenius structures using V1, V5, V6, and V8. In each case, the
analysis compares the unconstrained or soft reference, the two single-constraint
variants, and the joint hard-constraint variant.

\begin{table}[htbp]
    \centering
    \begin{tabular}{@{}llrrrl@{}}
        \toprule
        \textbf{Family}       & \textbf{Tier}   & \textbf{A only}
                              & \textbf{B only} & \textbf{Joint}  & \textbf{Interaction}                         \\
        \midrule
        Bounds + monotonicity & T1              & 5.624           & 5.213                & 7.116 & antagonistic  \\
        Bounds + monotonicity & T2              & 5.903           & 5.709                & 6.145 & antagonistic  \\
        Bounds + monotonicity & T3              & 2.217           & 4.643                & 1.192 & superadditive \\
        Sips + Van't Hoff     & T1              & 5.410           & 6.082                & 6.209 & antagonistic  \\
        Sips + Van't Hoff     & T2              & 5.476           & 5.543                & 5.056 & superadditive \\
        Sips + Van't Hoff     & T3              & 1.372           & 5.478                & 1.435 & superadditive \\
        \bottomrule
    \end{tabular}
    \caption{Full-protocol ablation interaction summary using RMSE values. The
        interaction label is computed from the factorial contrast in the ablation
        script.}
    \label{tab:ablation-interactions}
\end{table}

The full-protocol ablation results indicate that single constraints and joint
constraints do not behave identically across the evaluation tiers. For T1, the
single hard monotonicity model has the lowest RMSE, whereas the V2/V8 all-hard
combined variants give the lowest T2 RMSE. Therefore, the constraint associated
with the lowest error depends on whether the model is being tested in T1
interpolation or under coal-rank extrapolation. The interaction results are
also tier-dependent:
bounds and monotonicity are antagonistic on T1 and T2 but superadditive on T3.
The Sips and Van't Hoff combination is antagonistic on T1 and superadditive on
T2 and T3. The T3 findings remain uncertain because T3 contains only three
samples.

The pairwise statistical analysis gives the same warning. In the full-protocol
ablation results, 2 of 28 pairwise RMSE comparisons were significant after
Bonferroni correction on T1, and 7 of 28 were significant on T2. No pairwise
RMSE comparisons were significant on T3 after correction. This supports using
T3 mainly as descriptive pressure-extrapolation evidence rather than as strong
statistical evidence.

\subsection{Uncertainty calibration}

Uncertainty calibration was evaluated using 100 Monte Carlo dropout samples per
model and split. For each model, a scalar calibration factor \(\tau\) was fitted
on the validation set. This factor rescales the predicted standard deviation:
\(\tau=1\) leaves the uncertainty unchanged, while \(\tau>1\) makes the
prediction intervals wider. Calibration quality was measured using expected
calibration error (ECE), where smaller values indicate that the empirical
coverage of prediction intervals is closer to the nominal confidence levels.

\begin{table}[htbp]
    \centering
    \small
    \begin{tabular}{lrrrrr}
        \toprule
        \textbf{Variant}                      & \(\boldsymbol{\tau}\) & \textbf{Val. ECE}
                                              & \textbf{T1 ECE}       & \textbf{T2 ECE}   & \textbf{T3 ECE}                          \\
        \midrule
        V1 all-soft                           & 1.059                 & 0.032             & 0.027           & 0.273          & 0.120 \\
        V2 all-hard                           & 1.122                 & 0.026             & \textbf{0.015}  & \textbf{0.202} & 0.120 \\
        V3 hard bounds                        & 1.334                 & 0.042             & 0.053           & 0.242          & 0.120 \\
        V4 hard monotonicity                  & 0.891                 & 0.039             & 0.026           & 0.321          & 0.120 \\
        V5 hard Sips                          & 1.122                 & 0.041             & 0.032           & 0.276          & 0.120 \\
        V6 hard Van't Hoff / Arrhenius        & 1.259                 & 0.045             & 0.036           & 0.308          & 0.120 \\
        V7 hard bounds + monotonicity         & 1.000                 & \textbf{0.017}    & 0.015           & 0.299          & 0.120 \\
        V8 hard Sips + Van't Hoff / Arrhenius & 1.122                 & 0.026             & \textbf{0.015}  & \textbf{0.202} & 0.120 \\
        \bottomrule
    \end{tabular}
    \caption{Full-protocol calibrated ECE values for validation, T1 interpolation,
        T2 coal-rank extrapolation, and T3 pressure extrapolation. Lower ECE
        indicates better uncertainty calibration.}
    \label{tab:uq-calibration}
\end{table}

The full-protocol results show that validation calibration adjusts the uncertainty
estimates for most variants. The V7 hard-bounds and monotonicity model obtains
the lowest validation ECE, while the V2/V8 all-hard combined variants obtain
the lowest T1 and T2 ECE values in Table~\ref{tab:uq-calibration}.

The calibration results show that calibration on the validation set does not
always transfer to extrapolation. All variants have higher ECE on T2 than on
T1. This means that their uncertainty intervals are better calibrated for T1
interpolation samples than for the withheld
coal-rank extrapolation tier. In the full-protocol results, the V2/V8 all-hard combined
variants give the smallest T2 ECE, indicating that combining the hard Sips and
hard Van't Hoff / Arrhenius structures gives the lowest calibration error under
coal-rank extrapolation.

The epistemic uncertainty fraction also varies between variants. In general,
larger uncertainty under the held-out tiers is consistent with the idea that
the model expresses more model-form uncertainty when predictions are made
farther from the training distribution. However, because T3 contains only three
samples, the T3 uncertainty results should be treated as descriptive evidence
only.

\subsection{Interpretability and physical validation}

Interpretability analysis was used to check whether the trained models relied
on physically interpretable variables and whether the learned physical
parameters varied systematically across coal ranks. The analysis included SHAP-style feature
importance, learned Sips-parameter summaries, comparison with classical
Langmuir fits, and adsorption enthalpy summaries from the Van't Hoff /
Arrhenius model.

\begin{table}[htbp]
    \centering
    \begin{tabular}{rll}
        \toprule
        \textbf{Rank} & \textbf{V1 top feature} & \textbf{V5 top feature} \\
        \midrule
        1             & Pr                      & pressure\_MPa           \\
        2             & FC                      & FC                      \\
        3             & PT                      & volatile\_wt            \\
        4             & MV                      & MV                      \\
        5             & pressure\_MPa           & Tr                      \\
        \bottomrule
    \end{tabular}
    \caption{Top five global feature-importance rankings for the soft baseline
        and hard Sips-gauge model from the full-protocol interpretability analyses.}
    \label{tab:interpretability-shap}
\end{table}

For the hard Sips-gauge model, pressure was the most important input feature,
followed by fixed carbon and volatile matter. This is physically sensible
because the Sips-gauge architecture computes adsorption as a function of raw
pressure and learned coal-dependent Sips parameters. The V1 and V5 feature
rankings were weakly and not significantly aligned in the full-protocol analyses
(\(\rho_s=0.231\), \(p=0.471\)), showing that the hard Sips structure changes
how the model distributes importance across pressure, temperature, and
composition features.

\begin{table}[htbp]
    \centering
    \small
    \begin{tabular}{@{}llrrr@{}}
        \toprule
        \textbf{Split} & \textbf{Coal rank}         & \(\boldsymbol{\hat{q}_{\max}}\)
                       & \(\boldsymbol{\hat{n}}\)   & \(\boldsymbol{\Delta H}\)                       \\
        \midrule
        Train          & Anthracite                 & 24.66                           & 1.53 & -45.52 \\
        Train          & Low-volatile bituminous    & 20.80                           & 1.22 & -46.00 \\
        Train          & Medium-volatile bituminous & 15.70                           & 0.98 & -35.54 \\
        T2             & Subbituminous/lignite      & 10.25                           & 0.75 & -25.09 \\
        \bottomrule
    \end{tabular}
    \caption{Selected coal-rank summaries of learned physical quantities.
        \(\hat{q}_{\max}\) and \(\hat{n}\) are from V5, while
        \(\Delta H\) is from V6 and is reported in
        \(\mathrm{kJ\,mol^{-1}}\).}
    \label{tab:learned-physical-parameters}
\end{table}

The learned Sips parameters show a clear coal-rank pattern. The mean
\(\hat{q}_{\max}\) is highest for anthracite and lower for lower-rank coal.
Across all 114 experiments, \(\hat{q}_{\max}\) was negatively correlated with
volatile matter (\(r=-0.687\), \(p=3.04\times10^{-17}\)), satisfying the
blueprint check that higher volatile matter should generally correspond to
lower inferred adsorption capacity. The comparison with classical Langmuir
fits also supported this interpretation: all 114 experiment-level Langmuir
fits succeeded, and V5 \(\hat{q}_{\max}\) correlated with the fitted Langmuir
capacity \(q_L\) (\(r=0.739\), \(p=6.11\times10^{-21}\)).

The Van't Hoff / Arrhenius model produced negative adsorption enthalpy values,
consistent with exothermic adsorption. Mean \(\Delta H\) values were inside the expected
reference range used by the verification script for all reported ranks. The
rank effect was statistically significant for the training split, validation
split, and T1 split. For T2 and T3, rank-wise statistical comparison was not
possible because each contained only one inferred rank in the full-protocol split
definition.

Overall, the interpretability results support the physical reading of the hard
models: pressure controls the adsorption curve, coal composition affects the
learned capacity, and the inferred thermodynamic quantities vary with coal
rank. The interpretability analysis shows that the hard-constraint models can
be examined through physical quantities such as Sips parameters,
adsorption enthalpy, and coal-rank trends.

\subsection{Summary of findings}

The full-protocol analysis establishes the full evaluation structure for the study.
The dataset was split at the experiment level, with separate tiers for
T1 interpolation, T2 coal-rank extrapolation, and T3 pressure extrapolation.
This makes it possible to evaluate not only whether a model fits known
adsorption behavior, but also whether its predictions remain physically valid
when the test data move away from the training distribution.

Different hard-constrained models achieve the lowest RMSE on the three
evaluation tiers, but no single variant is best on every tier. The hard
monotonicity model gives the lowest T1 RMSE, the V2/V8 all-hard combined
variants give the lowest T2 RMSE, and the joint hard-bounds and monotonicity
model gives the lowest T3 RMSE. The
T3 result is descriptive rather than statistically conclusive because the T3 set
contains three samples. The
physical-verification checks show that the hard architectures enforce their
targeted structural properties, including non-negative adsorption, saturation
bounds, and pressure monotonicity.

The ablation analysis shows that physical constraints do not contribute in a
single uniform way. The lowest-RMSE constraint choice changes between T1, T2,
and T3. The uncertainty analysis shows that
calibration fitted on the validation set does not always transfer cleanly to
the coal-rank extrapolation tier. The interpretability analysis supports the
physical reading of the models: pressure and fixed carbon appear among the
leading features, learned Sips capacities vary with coal rank, and the
Van't Hoff / Arrhenius model
produces negative adsorption enthalpy values consistent with exothermic
adsorption.

The main finding is that hard constraints improve physical validity by
construction, while their effect on predictive accuracy and calibration depends
on the evaluation tier. The full V1 tuning protocol strengthens the comparison
with the soft baseline, but the T3 set contains only three samples and therefore
limits the strength of pressure-extrapolation claims.

\section{Conclusions}

The full-protocol results show that hard thermodynamic constraints provide a clear
advantage in physical validity, since the constrained architectures remove
specific classes of violations by construction. However, this advantage does
not translate into uniformly better predictive accuracy or uncertainty
calibration. The best-performing model depends on the evaluation tier: the hard
monotonicity model performs best on T1, the V2/V8 all-hard combined variants perform best on T2,
and the joint bounds--monotonicity model performs best on T3, although T3 is
too small for strong statistical claims. Thus, hard constraints outperform soft
penalties most clearly in guaranteed physical validity, while their accuracy
and calibration benefits are conditional.

\section{Further Work}

A fair comparison of computational cost remains to be investigated. The present
report compares predictive accuracy, physical validity, uncertainty
calibration, and interpretability, but it does not systematically investigate the computational cost.


\end{document}
